% Author: Loris Reiff
% Inspired by http://www.cadmo.ethz.ch/education/lectures/HS16/DA/uebungen/Loesung01_englisch.pdf
\documentclass[fleqn]{scrartcl}
\usepackage[utf8]{inputenc}
\usepackage[UKenglish]{babel}
\usepackage[T1]{fontenc}
\usepackage{lmodern}
\usepackage{amsmath,amssymb,amstext, amsthm}
\usepackage[headsepline,
  ]{scrlayer-scrpage}
\usepackage{tabto}
\usepackage{graphicx}
\usepackage{geometry}


%%%%%%%%%%%%%%%%%%%%%%%%%%%%%%%%%%%%%%%%%%%%%%%%%%%%%%%%%%%%%%%%%%%%%%%%
% Macros
\author{John Doe}									% Name
\date{\today}										% Date
\newcommand{\assignment}{1}							% Assignment Nr.
\newcommand{\class}{Mo. 9 -- 11, Room ABC 42}		% Exercise class
%%%%%%%%%%%%%%%%%%%%%%%%%%%%%%%%%%%%%%%%%%%%%%%%%%%%%%%%%%%%%%%%%%%%%%%%

% safe author and date
\makeatletter
\let\theauthor\@author
\let\thedate\@date
\makeatother


% page style
\geometry{a4paper, top=2.3cm, left=2.5cm, right=2.5cm} 
\clearpairofpagestyles % clear head and foot
\ihead{D\&A -- Assignment \assignment}
\chead{}
\ohead{\theauthor}
\addtokomafont{pagenumber}{\sffamily} % pagenumber font
\cfoot{\pagemark}
\setkomafont{pagehead}{\normalfont\sffamily}  % font for head
\setlength\parindent{0pt}	% no indent

% section style
\renewcommand{\thesubsection}{Exercise \thesection.\arabic{subsection}}
\setcounter{section}{\assignment}
\setkomafont{subsection}{\normalfont\itshape}
\makeatletter
\renewcommand*{\@seccntformat}[1]{% Section number formatting
   \protect\othersectionlevelsformat{#1}{%
     \expandafter\aftergroup\noexpand\@gobble}{\sffamily\bfseries\upshape\csname the#1\endcsname}%
}
\makeatother

% enumerate
\renewcommand{\labelenumi}{\alph{enumi})}

% ceil and floor
\newcommand{\ceil}[1]{\lceil #1 \rceil}
\newcommand{\floor}[1]{\lfloor #1 \rfloor}


\begin{document}
\begin{titlepage}
\hrule
\vspace{2mm}
\begin{minipage}{0.2\textwidth}
	\includegraphics[width=3cm]{res/eth_logo.eps}
	\end{minipage}
	\begin{minipage}{0.4\textwidth}
	\sffamily\small 
	Eidgenössische \\
	Technische Hochschule \\
	Zürich
\end{minipage}
\begin{minipage}{0.4\textwidth}
	\sffamily\small 
	Ecole polytechnique fédérale de Zurich\\
	Politecnico federale di Zurigo \\
	Federal Institute of Technology at Zurich
\end{minipage}
	
\vspace{2mm}
\hrule
\vspace{4mm}
\theauthor\hfill \thedate

\vspace{5mm}
{\fontfamily{lmss}\bfseries\selectfont\LARGE Data Structures and Algorithms \hfill Assignment
\assignment}
\vspace{2mm}

\def\arraystretch{1.3}
\begin{tabular}{@{}p{0.2\textwidth} p{0.75\textwidth}}
\Large{Exercise class:} & \class \\
\Large{Submitted by:} & \theauthor \\
\Large{Corrected by:} & \hrulefill \\
\Large{Bonus points:} & \hrulefill
\end{tabular}

\let\endtitlepage\relax
\end{titlepage}

% Begin of document
\subsection{Examples}
\begin{enumerate}
\item {
	We have $5n \leq n^2$ for $n \geq 5$, and $10 \leq n^2$ for $n \geq 4$ (in particular for $n \geq 5$).
	We choose $n_0 = 5$ and $c = 5$ and get
	\begin{align}
		3n^2 + 5n + 10 \leq 3n^2 + n^2 + n^2 = 5n^2 = cn^2 \text{ for all } n \geq n_0 = 5
	\end{align}
	Consequently, $3n^2 + 5n + 10 \in \mathcal{O}(n^2)$
}
\item {
	Suppose that $n^{\frac{1}{2}} \in \mathcal{O}(n^{\frac{1}{3}})$ Then there exist constants $c > 0$
	\dots
}
\end{enumerate}
\dots

\subsection{Simplifying Expressions in \ensuremath{\mathcal{O}} Notation.}

\begin{enumerate}
\item {$\mathcal{O}(2n + 14n^2) = \mathcal{O}(n^2)$}
\end{enumerate}
\dots
\end{document}
